\documentclass[11pt]{article}
\usepackage[letterpaper,margin=0.85in]{geometry}\usepackage[T1]{fontenc}\usepackage[utf8]{inputenc}\usepackage{amsmath,amssymb,enumitem,booktabs,array,xcolor,fancyhdr,tabularx}
\setlist{itemsep=2pt,topsep=3pt}\pagestyle{fancy}\fancyhf{}\renewcommand{\headrulewidth}{0.3pt}
\fancyhead[L]{\small 4.7 Part 1 $\cdot$ L'H?pital's Rule: Spot 0/0 and Use It Once}\fancyhead[R]{\small Make-up key}\fancyfoot[C]{\small\thepage}
\definecolor{ans}{RGB}{28,125,74}\newcommand{\A}[1]{\textcolor{ans}{\textbf{#1}}}
\begin{document}

\section*{4.7 Part 1 $\cdot$ L'H?pital's Rule: Spot 0/0 and Use It Once}
\textbf{You missed this lesson. Copy this into your notes. That is all you need to do.}\par\medskip
\subsection*{Question of the day}
What do I do when plugging in gives 0/0?\par
\subsection*{Targets}\begin{itemize}\item I can tell when plugging in gives 0/0, ?/?, a number, or a number over 0.
\item I can use L'H?pital's Rule once, taking the top and bottom derivatives separately.
\item I can say when the rule does not apply.\end{itemize}
\subsection*{Concepts}\begin{itemize}
\item \textbf{Plug in first} (plug in): Put $x=a$ into the top and bottom. A plain number means you are done.
\item \textbf{Indeterminate form} (0/0): $\frac00$ and $\frac{\infty}{\infty}$ tell you nothing yet. They need another way.
\item \textbf{L'H?pital's Rule} (lim f/g = lim f?/g?): For $\frac00$ or $\frac{\infty}{\infty}$: $\lim\dfrac fg=\lim\dfrac{f'}{g'}$. Differentiate top and bottom SEPARATELY.
\end{itemize}
\subsection*{What the class predicted, and what is true}\begin{itemize}
\item Near $x=1$, $f(x)=6(x-1)$ and $g(x)=3(x-1)$ both go to 0. The ratio $f(x)/g(x)$ goes to... \A{2.} The factors $(x-1)$ cancel, leaving $6/3=2$. The slopes 6 and 3 give the ratio.
\item $\displaystyle\lim_{x\to0}\dfrac{\sin x}{x}$ plugs in to $\dfrac00$. The limit is... \A{1.} Near 0, $\sin x$ looks like $x$, so the ratio is about 1. The tutorial shows the rule that proves it.
\end{itemize}
\section*{The tutorial, written out}
\subsection*{Concept tutorial (8 pictures)}
\textit{I can say what each plug-in result means for a limit.}\par
\begin{enumerate}
\item \textbf{Plug in first.} To find a limit as $x\to a$, first put $x=a$ into the top and the bottom. \textit{(Picture: forms. The table has three rows: what plugging in gives, what it means, and what to do next. The gold row is the one we are on.)}\\ \textit{Check:} Plug $x=2$ into $\dfrac{x+4}{x+1}$. What number is the limit? \A{2.}
\item \textbf{The form 0/0.} Plugging $x=2$ into $\dfrac{x^2-4}{x-2}$ gives $\dfrac00$. That tells us nothing yet. \textit{(Picture: forms. The second row is the form $\frac00$.)}\\ \textit{Check:} What does $\dfrac00$ tell us about the limit? \A{nothing yet: it needs another method.}
\item \textbf{Watch a ratio settle.} Both go to 0, but their ratio does not: slide $x$ toward 2 and watch $f\div g$. \textit{(Picture: ratio. Blue is $f(x)=x^2-4$ and orange is $g(x)=x-2$. Both reach 0 at $x=2$. The table lists $f\div g$ as $x$ moves toward 2. Slide $x$ along the steps.)}\\ \textit{Check:} Slide to step 5 ($x=2.01$). The ratio $f\div g$ is close to what? \A{4.}
\item \textbf{Slopes decide the ratio.} Near the point each function looks like a line. The ratio of the lines is the ratio of their slopes. \textit{(Picture: ratio. Two straight lines through $x=1$: $f$ has slope 6 and $g$ has slope 3. The ratio column is the same at every step.)}\\ \textit{Check:} $f$ has slope 6 and $g$ has slope 3. What does $f\div g$ settle at? \A{2.}
\item \textbf{Slopes are derivatives.} The slope of $f$ at 0 is $f'(0)=2$ and the slope of $g$ is $g'(0)=1$. So the ratio settles at $\dfrac{f'(0)}{g'(0)}=2$. \textit{(Picture: ratio. Blue is $f(x)=e^{2x}-1$ and orange is $g(x)=x$. Both are 0 at $x=0$. The ratio column settles near 2.)}\\ \textit{Check:} $f(x)=e^{2x}-1$ has $f'(0)=2$ and $g(x)=x$ has $g'(0)=1$. Find $\dfrac{f'(0)}{g'(0)}$. \A{2.}
\item \textbf{L'H?pital's Rule.} This is L'H?pital's Rule: for $\frac00$, the limit of $\dfrac fg$ equals the limit of $\dfrac{f'}{g'}$. \textit{(Picture: ratio. Blue is $\sin 3x$ and orange is $x$. The ratio settles at 3.)}\\ \textit{Check:} For $\dfrac{\sin3x}{x}$ at 0, $f'=3\cos3x$ and $g'=1$. Find $\dfrac{f'(0)}{g'(0)}$. \A{3.}
\item \textbf{Top and bottom SEPARATELY.} Differentiate the top and the bottom each on its own. Do NOT use the quotient rule on the whole fraction. \textit{(Picture: forms. The $\frac00$ row. The rule says to take the derivative of the top and of the bottom on their own.)}\\ \textit{Check:} For $\dfrac{x^2-4}{x-2}$, which is the L'H?pital step? \A{$\dfrac{2x}{1}$.}
\item \textbf{When the rule does NOT apply.} The rule works only for $\frac00$ or $\frac{\infty}{\infty}$. For $\dfrac30$ the limit is not a finite number, and the rule does not apply. \textit{(Picture: forms. The third row: a number over 0.)}\\ \textit{Check:} Plugging $x=1$ into $\dfrac{x+2}{x-1}$ gives $\dfrac30$. Can we use the rule? \A{No: it is not 0/0 or ?/?.}
\end{enumerate}
\subsection*{I do: the five moves}\begin{enumerate}\item Plug in
\item Name the form
\item Differentiate top and bottom
\item Plug in again\end{enumerate}
\subsection*{I do (7 steps)}
\textit{I can watch two 0/0 limits solved with the rule used once, and spot the common slip.}\par
\begin{enumerate}
\item \textbf{Plug in.} \textit{Move 1 $\cdot$ Plug in.} Put $x=0$ into the top and the bottom. \textit{(Picture: ratio. Blue is $e^{2x}-1$ and orange is $x$. At $x=0$ both are 0.)}\\ Top: $e^0-1=0$\\ Bottom: $0$\\ \textit{Check:} What is the top, $e^{2x}-1$, at $x=0$? \A{0.}
\item \textbf{Name the form.} \textit{Move 2 $\cdot$ Name the form.} Top 0 and bottom 0 is the form $\frac00$. The rule may be used. \textit{(Picture: forms. The gold row of the table is the form $\frac00$.)}\\ Form: $\dfrac00$\\ \textit{Check:} Which form is it? \A{$\dfrac00$.}
\item \textbf{Differentiate top and bottom.} \textit{Move 3 $\cdot$ Differentiate top and bottom.} Differentiate the top on its own and the bottom on its own. \textit{(Picture: ratio. The ratio of slopes: $f'=2e^{2x}$ and $g'=1$.)}\\ $\dfrac{(e^{2x}-1)'}{(x)'}=\dfrac{2e^{2x}}{1}$\\ \textit{Check:} Which line is the L'H?pital step? \A{$\dfrac{2e^{2x}}{1}$.}
\item \textbf{Plug in again.} \textit{Move 4 $\cdot$ Plug in again.} Now put $x=0$ into the new fraction. \textit{(Picture: ratio. The ratio is near 2 at the last step.)}\\ $\dfrac{2e^0}{1}=2$\\ So the limit is $2$\\ \textit{Check:} Find $\dfrac{2e^{0}}{1}$. \A{2.}
\item \textbf{Problem 2: a polynomial.} \textit{Moves 1 and 2 $\cdot$ Plug in and name.} \textbf{Problem:} Find $\displaystyle\lim_{x\to1}\dfrac{x^3-1}{x^2-1}$. Put $x=1$ into both: $1-1=0$ on top and $1-1=0$ on the bottom. \textit{(Picture: forms. The table row to check: is the form $\frac00$?)}\\ Top: $1^3-1=0$\\ Bottom: $1^2-1=0$\\ Form: $\dfrac00$\\ \textit{Check:} Which form does $x=1$ give? \A{$\dfrac00$.}
\item \textbf{Differentiate and plug in.} \textit{Moves 3 and 4.} The top becomes $3x^2$ and the bottom becomes $2x$. Then put $x=1$. \textit{(Picture: forms. The rule: top and bottom separately.)}\\ $\dfrac{3x^2}{2x}$\\ at $x=1$: $\dfrac32$\\ \textit{Check:} Find $\dfrac{3x^2}{2x}$ at $x=1$. \A{1.5.}
\item \textbf{The common slip.} Slip: using the quotient rule on the whole fraction. The rule takes the derivative of the top and of the bottom on their own. \textit{(Picture: forms. The quotient rule is for the derivative of a fraction. L'H?pital is a different tool.)}\\ Wrong: $\dfrac{3x^2(x^2-1)-(x^3-1)2x}{(x^2-1)^2}$\\ Right: $\dfrac{3x^2}{2x}$\\ \textit{Common slip.}\\ \textit{Check:} What is wrong with the first line? \A{It uses the quotient rule. The rule needs top and bottom derivatives separately..}
\end{enumerate}
\subsection*{Big ideas}\begin{itemize}\item Always plug in first. A plain number means you are done.
\item $\frac00$ and $\frac{\infty}{\infty}$ are indeterminate: the rule may be used.
\item L'H?pital: $\lim\dfrac fg=\lim\dfrac{f'}{g'}$, with top and bottom differentiated separately.
\item A number over 0 is not indeterminate. The rule does not apply.\end{itemize}
\section*{One worked sample of each skill}
\subsection*{Name the form}
\noindent\textbf{Level 1.} Find $\displaystyle\lim_{x\to1}\dfrac{x^2-1}{x-1}$. Put $x=1$ into the top and bottom. What happens?\\ \A{(A) It gives $\frac00$, so L'H?pital's Rule may be used.} \quad (B) It gives a number over 0, so the rule does not apply. \quad (C) It gives a plain number, so just plug in. \quad (D) It gives $\frac{\infty}{\infty}$, so L'H?pital's Rule may be used.\\ \A{Answer: A} \quad Both the top and the bottom go to 0. That form tells us nothing yet. It gives $\frac00$, so L'H?pital's Rule may be used.\par\medskip
\noindent\textbf{Level 2.} Find $\displaystyle\lim_{x\to0}\dfrac{e^x-1}{x}$. Put $x=0$ into the top and bottom. What happens?\\ (A) It gives $\frac{\infty}{\infty}$, so L'H?pital's Rule may be used. \quad \A{(B) It gives $\frac00$, so L'H?pital's Rule may be used.} \quad (C) It gives a plain number, so just plug in. \quad (D) It gives a number over 0, so the rule does not apply.\\ \A{Answer: B} \quad Both the top and the bottom go to 0. That form tells us nothing yet. It gives $\frac00$, so L'H?pital's Rule may be used.\par\medskip
\noindent\textbf{Level 3.} Find $\displaystyle\lim_{x\to0}\dfrac{1-\cos x}{x^2}$. Put $x=0$ into the top and bottom. What happens?\\ (A) It gives $\frac{\infty}{\infty}$, so L'H?pital's Rule may be used. \quad (B) It gives a number over 0, so the rule does not apply. \quad (C) It gives a plain number, so just plug in. \quad \A{(D) It gives $\frac00$, so L'H?pital's Rule may be used.}\\ \A{Answer: D} \quad Both the top and the bottom go to 0. That form tells us nothing yet. It gives $\frac00$, so L'H?pital's Rule may be used.\par\medskip
\subsection*{The L'H?pital step}
\noindent\textbf{Level 1.} $\dfrac{x^2-9}{x-3}$ gives $\frac00$. Which line is the L'H?pital step?\\ \A{(A) $\dfrac{2x}{1}$} \quad (B) $\dfrac{x^2-9}{1}$ \quad (C) $\dfrac{2x}{x-3}$ \quad (D) $\dfrac{(2x)(x-3)-(x^2-9)(1)}{(x-3)^2}$\\ \A{Answer: A} \quad Top: $(x^2-9)'=2x$. Bottom: $(x-3)'=1$. The step is $\dfrac{2x}{1}$.\par\medskip
\noindent\textbf{Level 2.} $\dfrac{e^x-1}{x}$ gives $\frac00$. Which line is the L'H?pital step?\\ (A) $\dfrac{(e^x)(x)-(e^x-1)(1)}{(x)^2}$ \quad \A{(B) $\dfrac{e^x}{1}$} \quad (C) $\dfrac{e^x}{x}$ \quad (D) $\dfrac{e^x-1}{1}$\\ \A{Answer: B} \quad Top: $(e^x-1)'=e^x$. Bottom: $(x)'=1$. The step is $\dfrac{e^x}{1}$.\par\medskip
\noindent\textbf{Level 3.} $\dfrac{\sin 2x}{\sin 5x}$ gives $\frac00$. Which line is the L'H?pital step?\\ (A) $\dfrac{(2\cos 2x)(\sin 5x)-(\sin 2x)(5\cos 5x)}{(\sin 5x)^2}$ \quad (B) $\dfrac{\sin 2x}{5\cos 5x}$ \quad (C) $\dfrac{2\cos 2x}{\sin 5x}$ \quad \A{(D) $\dfrac{2\cos 2x}{5\cos 5x}$}\\ \A{Answer: D} \quad Top: $(\sin 2x)'=2\cos 2x$. Bottom: $(\sin 5x)'=5\cos 5x$. The step is $\dfrac{2\cos 2x}{5\cos 5x}$.\par\medskip
\subsection*{Find the limit}
\noindent\textbf{Level 1.} Find $\displaystyle\lim_{x\to3}\dfrac{x^2-9}{x-3}$. Give it to 2 decimal places if needed.\\ \A{Answer: 6} \quad Both go to 0, so use the rule. $\dfrac{2x}{1}$ at $x=3$ gives $6$.\par\medskip
\noindent\textbf{Level 2.} Find $\displaystyle\lim_{x\to0}\dfrac{e^{6x}-1}{x}$. Give it to 2 decimal places if needed.\\ \A{Answer: 6} \quad Both go to 0, so use the rule. $\dfrac{6e^{6x}}{1}$ at $x=0$ gives $6$.\par\medskip
\noindent\textbf{Level 3.} Find $\displaystyle\lim_{x\to1}\dfrac{x^2-1}{x^3-1}$. Give it to 2 decimal places if needed.\\ \A{Answer: 0.67} \quad Both go to 0, so use the rule. $\dfrac{2x}{3x^2}$ at $x=1$ gives $\dfrac{2}{3}$.\par\medskip
\subsection*{Spot the slip}
\noindent\textbf{Level 1.} A student uses L'H?pital's Rule on $\dfrac{x^2-9}{x-3}$ at $x=3$ and writes $\dfrac{2x(x-3)-(x^2-9)}{(x-3)^2}$. What is wrong?\\ \A{(A) It used the quotient rule instead of separate derivatives.} \quad (B) It used the right method but made an arithmetic error. \quad (C) There is no error. \quad (D) It should have differentiated the top twice.\\ \A{Answer: A} \quad It used the quotient rule instead of separate derivatives. L'H?pital takes the derivative of the top and of the bottom on their own: $\dfrac{2x}{1}$.\par\medskip
\noindent\textbf{Level 2.} A student uses L'H?pital's Rule on $\dfrac{x^2-36}{x-6}$ at $x=6$ and writes $\dfrac{2x(x-6)-(x^2-36)}{(x-6)^2}$. What is wrong?\\ (A) It used the right method but made an arithmetic error. \quad (B) It should have differentiated the top twice. \quad \A{(C) It used the quotient rule instead of separate derivatives.} \quad (D) There is no error.\\ \A{Answer: C} \quad It used the quotient rule instead of separate derivatives. L'H?pital takes the derivative of the top and of the bottom on their own: $\dfrac{2x}{1}$.\par\medskip
\noindent\textbf{Level 3.} A student writes: The limit of $\dfrac{x+4}{x+1}$ as $x\to2$: "this is a limit, so use the rule: $\dfrac11=1$." What is wrong?\\ (A) There is no error. \quad (B) It used the right method but made an arithmetic error. \quad (C) It should have differentiated the top twice. \quad \A{(D) The form was not $\frac00$, so the rule was not allowed. Plug in instead.}\\ \A{Answer: D} \quad The form was not $\frac00$, so the rule was not allowed. Plug in instead. At $x=2$ the top is 6 and the bottom is 3, a plain number. Just plug in: $\dfrac{6}{3}$.\par\medskip
\end{document}